\subsection{Predeclared real-energy decision study}
We evaluated the operational use of the certificate on the three electricity-consumption zones of the Tetouan City data set. The original ten-minute observations were aggregated into nonoverlapping two-hour means, giving 4,368 observations per zone. Weather variables were deliberately excluded so that the study remained a univariate forecasting comparison. Each zone was divided chronologically into 20\% prior, 45\% bound, 15\% validation, and 20\% test segments. The compared families were ridge regression, a fixed eight-rule dense TSK model, and a radius-controlled dense TSK model. Candidate lag orders were 3, 6, and 12, where lag 12 corresponds to 24 hours.

The prior, scaling, antecedent geometry, and clipping interval were constructed from the prior segment only. We used $B=\max\{1,1.25\max_{t\in D_{\rm prior}}|Y_t|\}$ on the prior-standardized scale. The certificate used the bound and validation segments, with total confidence $0.05$ divided familywise across the three zones. Before opening the test segment, a candidate was declared deployment-eligible only when its certificate did not exceed 0.20 and its certification target-clipping rate did not exceed 5\%. Among eligible candidates, the predeclared decision selected the one with the lowest validation RMSE, provided it was no worse than a 24-hour seasonal-naive fallback. Underforecast error received weight two and overforecast error weight one in the operational cost. All decisions were serialized before test forecasts were computed.
